\documentclass[11pt]{article}

\usepackage[left=0.8in,right=0.8in,top=0.9in,bottom=0.9in]{geometry}
\usepackage{amsmath}
\usepackage{enumitem}
\usepackage[hidelinks]{hyperref}
\usepackage{microtype}
\usepackage[backend=biber, style=numeric]{biblatex}


\linespread{0.97}

\setlength{\parindent}{15pt}
\setlength{\parskip}{4pt}

\hypersetup{
    colorlinks=true,
    urlcolor=blue
}

\addbibresource{../references.bib}

\setlist[itemize]{leftmargin=*, itemsep=1pt, topsep=2pt}
\begin{document}

\begin{center}
{\Large \textbf{Orchestrating Efficient LLM Inference via Structural and Adaptive
Speculative Decoding}}\\[7pt]
COMPSCI 690AB - Systems for Deep Learning\\[4pt]
Harish Srinivasan (hsrinivasan@umass.edu), Pooja Rajeev (prajeev@umass.edu) \\
\noindent\textbf{Repository:} \href{https://github.com/Sykarius/speculative-decoding}{GitHub - speculative-decoding}
\end{center}


\noindent\textbf{Current Status:} On Track / Phase 1 Completed (Baseline + Some SD strategies implemented)
\section{Summary}
We have successfully moved past the architectural baseline phase and completed the implementation of four primary inference strategies. 
The project has transitioned from pure development to the optimization and benchmarking phase. 
Crucially, initial sanity checks confirm that all implemented Speculative Decoding (SD) methods outperform the standard autoregressive baseline, validating the underlying logic of our acceleration framework.
The implementation including the benchmarking scripts can be found in the project repository: \href{https://github.com/Sykarius/speculative-decoding}{GitHub - speculative-decoding}

\section{Accomplishments}
\begin{itemize}
    \item \textbf{Baseline Established:} Completed the reference autoregressive approach to serve
    as the performance baseline for all SD variations
    \item \textbf{Core Speculative Logic:} Fully implemented the speculative decoding method
    \begin{itemize}
        \item \textbf{Fixed Window:} Implemented with both standard sampling \cite{leviathan2023} and stochastic rejection sampling \cite{leviathan2023}
        \item \textbf{Adaptive Lookahead:} Implemented a simple AIMD \cite{chiu1989analysis} (Additive Increase / Multiplicative Decrease) controller to adjust lookahead($\gamma$) based on acceptance outcomes
    \end{itemize}
    \item \textbf{Benchmarking Infrastructure:} Implemented a comprehensive benchmarking script to capture telemetry and log performance metrics to the disk. Currently it captures TTFT, TPOT, throughput (tokens/sec), acceptance rate ($\alpha$) and efficiency of the drafting
\end{itemize}

\section{Upcoming Milestones}
\begin{itemize}
    \item \textbf{Adaptive Strategies:} We will expand the adaptive strategies to include Entropy-based and Top-K adaptive windowing. (mentioned in AdaSD \cite{lu2025adasdadaptivespeculativedecoding})
    \item \textbf{Architectural Optimization:} Developing Block Verification \cite{sun2025blockverificationacceleratesspeculative} and Tree-Based Speculation \cite{Miao_2024} to handle multiple candidate sequences simultaneously
    \item \textbf{Large-Scale Validation:} Transitioning from sanity checks to a full-dataset benchmark and to generate final performance metrics. Perform the validation with different parameter ($\gamma$) values
    \begin{itemize}
        \item \textbf{Draft Model Sweep:} We run validation with different lookahead ($\gamma$) values to thoroughly understand the effects of draft size.
        \item \textbf{Comparative Study:} Compare the different SD algorithms across various metrics to gain a comprehensive understanding of their performance
    \end{itemize}
    \item \textbf{Final Analysis:} Organize and analyze the benchmark result and derive key performance insights
\end{itemize}

\nocite{*}

\printbibliography
\enlargethispage{1\baselineskip}
\end{document}